\section{Compile Time and Pipeline Stages}\label{sec:compile-time}

A Motivo compilation runs four sequential stages, parsing, semantic analysis, lowering, and MIDI serialization, and the
harness times each one independently (Section~\ref{sec:evaluation-methodology}).
This decomposition is the through-line of the evaluation: across every program measured in this chapter, parsing and
semantic analysis stay below $0.07$~ms (and below $0.16$~ms even for the most verbose real composition), while lowering
and MIDI serialization account for essentially all of the compiler's work.
The cost of a compilation is therefore governed almost entirely by the timeline it produces and by how much work each
emitted event demands.
This section establishes how that cost grows with the size of the output; Section~\ref{sec:scaling-stress} then shows
how it varies with the shape of the source program at a fixed output size.

\subsection{Scaling with Output Size}\label{detail-subsec:scaling-sweep}

The scaling sweep compiles a flat \texttt{loop (N)} that emits one note per iteration for
$N \in \{1000, 2000, 5000, 10000, 15000, 20000\}$.
Holding the source shape constant isolates the cost of materializing and serializing progressively larger timelines.
Table~\ref{tab:scaling-stages} reports the per-stage means, and Figure~\ref{fig:scaling-sweep} plots them against the
event count.

\begin{table}[htbp]
  \centering
  \footnotesize
  \renewcommand{\arraystretch}{1.25}
  \begin{tabular}{rrrrrr}
    \toprule
    \textbf{Note events} & \textbf{Parse} & \textbf{Semantic} & \textbf{Lowering} & \textbf{MIDI write} &
    \textbf{End-to-end} \\
    \midrule
    1,000 & 0.019 & 0.004 & 0.079 & 0.206 & $2.42 \pm 0.06$ \\
    2,000 & 0.019 & 0.004 & 0.156 & 0.275 & $2.56 \pm 0.06$ \\
    5,000 & 0.019 & 0.004 & 0.384 & 0.621 & $3.17 \pm 0.09$ \\
    10,000 & 0.020 & 0.004 & 0.765 & 1.281 & $4.25 \pm 0.09$ \\
    15,000 & 0.019 & 0.004 & 1.150 & 1.818 & $5.17 \pm 0.12$ \\
    20,000 & 0.020 & 0.004 & 1.525 & 2.549 & $6.30 \pm 0.12$ \\
    \bottomrule
  \end{tabular}
  \caption{Per-stage mean compile time (ms) for the scaling sweep, current \texttt{motivoc}. Parse and semantic times
    are flat; lowering and MIDI write grow with the event count. The end-to-end column reports the mean $\pm$ standard
  deviation over 50 runs and additionally includes process startup and file I/O.}\label{tab:scaling-stages}
\end{table}

\begin{figure}[htbp]
  \centering
  \begin{tikzpicture}
    \begin{axis}[
        width=0.92\textwidth,
        height=6.6cm,
        xlabel={Note events},
        ylabel={Mean time (ms)},
        xmin=0, xmax=21000,
        ymin=0,
        xtick={0,5000,10000,15000,20000},
        scaled x ticks=false,
        x tick label style={/pgf/number format/1000 sep={,}},
        legend pos=north west,
        legend style={font=\scriptsize, draw=gray!40},
        grid=major,
        grid style={gray!18},
        tick label style={font=\scriptsize},
        label style={font=\small},
      ]
      \addplot[mark=*, motivoBlue, very thick] coordinates {
      (1000,2.418) (2000,2.555) (5000,3.167) (10000,4.253) (15000,5.173) (20000,6.298)};
      \addplot[mark=square*, motivoCyan, very thick] coordinates {
      (1000,0.206) (2000,0.275) (5000,0.621) (10000,1.281) (15000,1.818) (20000,2.549)};
      \addplot[mark=triangle*, motivoAmber, very thick] coordinates {
      (1000,0.079) (2000,0.156) (5000,0.384) (10000,0.765) (15000,1.150) (20000,1.525)};
      \legend{End-to-end, MIDI write, Lowering}
    \end{axis}
  \end{tikzpicture}
  \caption{Compile time versus event count for the scaling sweep (current \texttt{motivoc}, Release).
    Lowering and MIDI write grow with the event count; the end-to-end curve sits about $2$~ms higher because it includes
  the fixed process-startup and file-I/O floor that the in-process stage timers exclude.}\label{fig:scaling-sweep}
\end{figure}

Two effects are visible.
First, the end-to-end curve does not pass through the origin: even the 1{,}000-event program takes about $2.4$~ms,
almost all of it process startup and filesystem overhead, because its lowering and MIDI stages together cost under
$0.3$~ms.
This fixed floor dominates small programs and explains why the real compositions in Section~\ref{sec:general-case}, none
of which exceeds a few hundred events, all compile in roughly the same time despite differing in size.
Second, once that floor is removed, the two heavy stages scale with the event count as their algorithms predict:
lowering grows close to linearly ($O(E)$) and MIDI serialization slightly faster, consistent with the $O(E \log E)$
tick-sort in the backend.
Doubling the events from 10{,}000 to 20{,}000 raises lowering by $2.00\times$ and MIDI write by $2.00\times$, and MIDI
write remains the larger of the two at every size because it both sorts and byte-encodes each event.

\subsection{Interactive Latency in the Studio}\label{detail-subsec:execution-speed-and-real-time-viability}

Motivo Studio compiles when the user presses the compile button, not on every keystroke or on save.
At roughly $5$--$8$~ms for typical and flat programs at the 20{,}000-event cap, and up to about $9$~ms for the most
expression-heavy shapes (Section~\ref{sec:scaling-stress}), that whole-file compile step is fast enough that the
editor's compile-to-preview loop feels immediate in practice.
Per-keystroke recompilation would require incremental parsing and partial lowering, which the current architecture does
not implement, but the measured throughput leaves ample headroom if such a feature were added later.
Real compositions of a few hundred events (Section~\ref{sec:general-case}) compile faster still, so auditory and visual
feedback never depends on an interpreted runtime at playback.

\subsection{The Cost of Explicit Typing}\label{detail-subsec:typing-cost}

Motivo version~2.0 replaced the implicit \texttt{let} typing of version~1.0 with explicit declarations and
signature-based overload resolution (Section~\ref{sec:semantic-rules}).
A natural question is whether those stricter static guarantees cost anything at compile time.
Three fixtures exist in both source dialects, a flat 20{,}000-event loop, a two-note chord loop of 9{,}980 events, and a
triple-nested loop cascade of 20{,}000 events, and Table~\ref{tab:typing-cost} compiles each with both builds.

\begin{table}[htbp]
  \centering
  \footnotesize
  \renewcommand{\arraystretch}{1.25}
  \begin{tabular}{llrrrrr}
    \toprule
    \textbf{Fixture} & \textbf{Ver.} & \textbf{Parse} & \textbf{Semantic} & \textbf{Lowering} & \textbf{MIDI write} &
    \textbf{End-to-end} \\
    \midrule
    Flat loop (20,000)        & v1 & 0.021 & 0.004 & 1.559 & 2.558 & $6.40 \pm 0.32$ \\
    & v2 & 0.019 & 0.004 & 1.552 & 2.569 & $6.37 \pm 0.32$ \\
    \midrule
    Chord loop (9,980)        & v1 & 0.021 & 0.004 & 0.858 & 1.261 & $4.32 \pm 0.12$ \\
    & v2 & 0.020 & 0.004 & 0.852 & 1.247 & $4.29 \pm 0.09$ \\
    \midrule
    Nested cascade (20,000)   & v1 & 0.022 & 0.004 & 2.056 & 2.572 & $6.91 \pm 0.14$ \\
    & v2 & 0.021 & 0.004 & 2.064 & 2.580 & $6.91 \pm 0.12$ \\
    \bottomrule
  \end{tabular}
  \caption{Per-stage mean compile time (ms) for the two compiler versions on the three dual-dialect fixtures; the
    end-to-end column reports the mean $\pm$ standard deviation over 50 runs. Each pair produced byte-identical MIDI
  (180{,}056\,/\,84{,}886\,/\,180{,}056~B).}\label{tab:typing-cost}
\end{table}

The two versions are indistinguishable.
End-to-end means differ by between $-0.7\%$ and $0.0\%$ across the three fixtures, far inside run-to-run noise: the
per-fixture standard deviations in Table~\ref{tab:typing-cost} are $0.1$--$0.3$~ms, an order of magnitude larger than
the sub-$0.05$~ms differences between versions.
Every fixture also produces byte-identical MIDI output, confirming that the type-system change did not alter what the
compiler emits.
The per-stage breakdown explains the result: parsing and semantic analysis are already sub-$0.025$~ms in both versions,
so the added type annotations and signature-based lookup are lost in the noise, while lowering and MIDI write, which
dominate the total, are untouched by the frontend change.
Collapsing version~1.0's two semantic traversals into one in version~2.0 yields no measurable speedup for the same
reason: both were already negligible.
The practical conclusion is that the move to explicit typing buys stronger static guarantees and a simpler semantic
architecture at no compile-time cost.
